\documentclass[a4paper,12pt]{article}

\usepackage[utf8]{inputenc}   % Umlaute
\usepackage[T1]{fontenc}
\usepackage[ngerman]{babel}   % Deutsche Sprache

\title{Dynamic Time Warping}
\author{Max Mustermann}
\date{\today}

\begin{document}

\maketitle

\section{Einleitung}
Dies ist ein kurzes Beispiel für ein LaTeX-Dokument.  
Man kann damit wissenschaftliche Texte, Berichte und Bücher setzen.

\section{Beispiel}
Hier ein kleines Beispiel für eine einfache Formel:
\[
E = mc^2
\]

\section{Fazit}
LaTeX ist sehr hilfreich, wenn man strukturierte Dokumente erstellen möchte.
\cite{einstein1920}
\bibliographystyle{plain}
\bibliography{late.bib}

\end{document}
